\clearpage
\phantomsection\addcontentsline{toc}{section}{Glossario}
\markboth{Glossario}{Glossario}
{\headfont\bfseries\Huge\color{navy}Glossario\par}
\vspace{4pt}{\color{teal}\rule{3cm}{2.5pt}}\vspace{8pt}

\newcommand{\voce}[2]{\par\vspace{4pt}{\headfont\bfseries\small\color{navy}#1}\enspace #2\par}
\voce{Attributo}{Proprietà elementare di un'entità o di una relationship; associa a ogni occorrenza un valore di un dominio.}
\voce{Cardinalità}{Coppia $(\min,\max)$ di partecipazioni di un'occorrenza di entità a una relationship (o di valori di un attributo).}
\voce{Dizionario dei dati}{Documentazione testuale di entità e relationship di uno schema.}
\voce{Entità}{Classe di oggetti della realtà con proprietà comuni ed esistenza autonoma.}
\voce{Entità debole}{Nei testi in notazione Chen: entità identificata attraverso un'entità proprietaria; corrisponde all'identificatore esterno.}
\voce{Ereditarietà}{Le figlie di una generalizzazione hanno tutte le proprietà del genitore.}
\voce{Generalizzazione}{Legame fra un'entità genitore e figlie che ne sono casi particolari; totale/parziale, esclusiva/sovrapposta.}
\voce{Identificatore}{Insieme minimale di attributi (interno) o di attributi ed entità esterne (esterno) che distingue le occorrenze di un'entità.}
\voce{Occorrenza}{Un singolo elemento di un'entità (l'oggetto stesso) o di una relationship (una coppia o $n$-upla).}
\voce{Relationship}{Legame logico fra due o più entità; le sue occorrenze non si ripetono.}
\voce{Ruolo}{Nome del «posto» di un'entità in una relationship ricorsiva.}
\voce{Schema concettuale}{Descrizione della realtà di interesse indipendente dal DBMS, nel modello E-R.}
\voce{Sottoinsieme}{Generalizzazione con una sola entità figlia.}
\voce{Vincolo non esprimibile}{Regola della realtà che i costrutti del modello non riescono a rappresentare; si documenta a parole.}
